\section{Data files}
\label{sec:data}

Every model needs data in order to proceed to its calibration, so the
possibilty to enter data in a convenient format is important.
\amlet{} unfortunately currently supports basic formats only, even if
development is planned on this issue.
The supported formats reflect the original purpose of the package,
which was only testing a new optimisation method.

\subsection{\gauss{}-like formats}

Preliminary \amlet{} design mimics Train's \gauss{} codes for mixed
logit estimation (cross-sectional and panel versions).
These routines are available at the address
\url{http://elsa.berkeley.edu/~train/software.html}.
We also refer the reader to Revelt and Train~\cite{ReveTrai98}, and
Train~\cite{Trai99}, for more details.
\amlet{} has currently no capabilities to construct new data from a
preexisting dataset, for instance  dummy
observations for alternatives specific constants, so please be sure
that your data are completely specified before using \amlet{}.

The format of the data file is specified by assign to the variable
\verb|format| of the driver file the value \verb|Train96| of
\verb|Train99|, the number being a reference to Train's code release
year.
For instance, the following line
\begin{verbatim}
format = Train96;
\end{verbatim}
indicated that data are stored in \verb|Train96| format, which is also
taken as default if no format is specified in the driver file.

\subsubsection{Train96 format}

This format is primarily designed for cross-sectional data, as
illustrated in Figure~\ref{fig:Train96}.
In this figure, each indiviual is facing four alternatives,
corresponding four columns.
The six first rows of each record (containing seven rows) corresponds
to the independent variables.
Since each individual is assumed to deliver one observation, this
format cannot be used for panel data.
An additional row is devoted to the dependent (decision) variable,
with one on the column corresponding  to the chosen alternatives, and
0 elsewhere.
The line position in each record must be indicated with the variable
\verb|iddep| in the driver file.
For Figure~\ref{fig:Train96}, we would have
\begin{verbatim}
iddep = 7;
\end{verbatim}
\begin{figure}
\begin{verbatim}
     0    1    0    0
     0    0    1    0
     0    0    0    1
   1.6 -2.0  3.2  4.1
     0  2.5  1.0  0.8
  -1.2 -1.5 -0.4 -1.6
     0    1    0    0
     0    1    0    0
     0    0    1    0
     0    0    0    1
   2.4 -3.0  2.3  3.0
     0  1.7  1.2  1.5
  -1.6 -2.0 -1.2 -2.4
     0    0    0    1
\end{verbatim}
\label{fig:Train96}
\caption{Train96 Format}
\end{figure}
%\textit{Cas o� toutes les alternatives ne sont pas disponibles}

By default, all alternatives are available to each individual.
If you want to specify alternative availabilities, you have to set the
censor variable to 1 in the driver file:
\begin{verbatim}
censor = 1;
\end{verbatim}
The variable \verb|idcensor| then represents the line describing
alternative availabilities.
The corresponding line in the driver file is simply a 0-1 line, with
one on available alternatives, and 0 elsewhere.

\subsubsection{Train99 format}

This format is used for panel data.
It adds to the base name the suffixes \verb|_x|, \verb|_y| and \verb|_times|.
